After working on DreamDiffusion, it was asked to get the bigger picture to
compare DreamDiffusion to more recent papers in the same domain.

The generated images from DreamDiffusion showed unpleasant results as they
failed to capture the image structure or the semantic of the image's content.
For example, if the ground truth image is a dog in the middle of a football
field by a sunny day, the corresponding generated image can be a dog in a flat
eating snacks. This problem gave us the idea of making a paper about evaluating
the EEG to image generators on their semantic faithfulness.

A first state of the art of EEG to image generators had to be done, we had to
figure out which of these were reproducible so we could evaluate their semantic
performance \cref{fig:stateoftheart}, very few were reproducible, quite often
the code, the dataset or the checkpoints were not available and the researchers
took too much time to send them.

\insertfigure
    {assets/images/stateoftheart}
    {width=\textwidth}
    {State of the art of EEG to image generators}
    {fig:stateoftheart}

\subsubsection{Semantic faithfulness metric}

It was decided to focus only on diffusion models only, and the semantic
faithfulness metric was to be found. Few papers addressed this problem, the
NECOMIMI \cite{necomimi} paper addressed this with their own metric,
the Category-based Assessment Table (CAT) score. The generated and ground truth
images were fed to ChatGPT 4o with the following prompt: "Please provide me with
5 one-word descriptions of the image, ranging from high level to low level"
resulting in 5 words, if two words matched it gave +1 to the score, resulting in
a score from 0 to 5 for one generated image, this score was then averaged to
evaluate models on the test set.

DreamSim \cite{dreamsim} is a very interesting paper where they focus on
"perceptual similarity" between the ground truth and the generated image. It
means that the score is matching what an human would describe as two likely
images from its own criteria. This perceptual similarity is trained from
different existing scores that are going through an MLP that gives a final
score. This model is trained with a dataset of images where real humans decided
wether image A or C matched the best the image B. The model is made to match the
humans decision. (IMAGE FROM \cite{dreamsim})\cref{fig:dreamsim-collection}
shows what was shown to the humans to make the dataset.

\insertfigure
    {assets/images/dreamsim-collection}
    {width=0.4\textwidth}
    {Data collection for DreamSim (FROM \cite{dreamsim})}
    {fig:dreamsim-collection}

\subsubsection{Another image generator, BrainVis}

\insertfigure
    {assets/images/brainvis-good}
    {width=0.7\textwidth}
    {Good images generated by BrainVis (ground truth on the left)}
    {fig:brainvis-good}

The only other model easy to reproduce model found was BrainVis \cite{brainvis}.
Compared to DreamDiffusion it generates very few incoherent images, the most
incoherent ones seem to be with faces, digits, or parachutes like
\cref{fig:brainvis-incoherent}

\insertfigure
    {assets/images/brainvis-incoherent}
    {width=0.7\textwidth}
    {Almost incoherent images generated by BrainVis (ground truth on the left)}
    {fig:brainvis-incoherent}

\insertfigure
    {assets/images/brainvis-bad}
    {width=0.7\textwidth}
    {Bad image class generated by BrainVis (ground truth on the left)}
    {fig:brainvis-bad}

BrainVis image generation was much more easier to conduct as the code and the
instructions were clear.

\subsubsection{Experiments to conduct}

To have a score that captures the semantics of the images we thought it was
relevent to first go from image to text through a captioning model like BLIP2
\cite{blip2} or MiniCPM \cite{minicpm}.

The experiments in \cref{fig:experiments} were planned to be conducted but no
third captioning model was found to be tested and there was not enough time
during the internship to conduct them all.

\insertfigure
    {assets/images/experiments}
    {width=\textwidth}
    {Experiments planned at the end of the internship}
    {fig:experiments}

The idea of the experiments in \cref{fig:experiments} was to have four sets of
generated images, two by dataset and two by model. Then, for each (ground truth,
generated) image pair, use a captioning model (giving a text description of the
image), that we can then feed to a sentence encoder model (sBERT \cite{sbert}),
we would end up with a pair of two sBERT embeddings we would apply our cosine
similarity to to have our score. By first extracting the caption of the image we
would be more semantically correct on the score. \cref{fig:experiments-pipeline}
shows the whole pipeline.

\insertfigure
    {assets/images/experiments-pipeline}
    {width=\textwidth}
    {Experiments pipeline for semantically accurate score}
    {fig:experiments-pipeline}

The only experiments that could be conducted were on the CVPR40 dataset with
DreamDiffusion and BrainVis, computing the following image similarity metrics:
PSNR\cite{psnr} score, SSIM\cite{ssim} score, LPIPS\cite{lpips} score,
CLIP\cite{clip} score, DISTS\cite{dists} score, DINO\cite{dino} score,
DreamSim\cite{dreamsim} score and our BLIP2 captioning score. The results are in
\cref{table:exps-results}.

\begin{table}[h!]
    \centering
    \label{table:exps-results}
    \caption{Evaluations of BrainVis and DreamSim on different scores, including
    ours}
    \begin{tabular}{lcc}
    \hline
    \textbf{Metric} & \textbf{BrainVis} & \textbf{DreamDiffusion} \\
    \hline
    PSNR & 56.79 & 57.25 \\
    SSIM  & 0.189 & 0.174 \\
    LPIPS & 0.818 & 0.751 \\
    CLIP  & 0.662 & 0.544 \\
    DISTS & 0.397 & 0.394 \\
    DINO  & 0.306 & 0.082 \\
    DreamSim & 0.666 & 0.775 \\
    Ours & 0.410 & 0.188 \\
    \hline
    \end{tabular}
\end{table}

Having not enough time to conduct these experiments it was not possible to
submit any paper by the end of the internship, however the project was continued
and I volunteered to plan the experiments for a paper to be submitted in a
conference.